%************************************************************************
% Gemeinsame Preamble fuer Skript (docu.tex), Folien (slides.tex) und
% Handouts/Uebungsblaetter (handouts.tex).
% Jede Root-Datei laedt ihre Dokumentklasse, setzt ggf. \slidestrue bzw.
% \handoutstrue, dann %************************************************************************
% Gemeinsame Preamble fuer Skript (docu.tex), Folien (slides.tex) und
% Handouts/Uebungsblaetter (handouts.tex).
% Jede Root-Datei laedt ihre Dokumentklasse, setzt ggf. \slidestrue bzw.
% \handoutstrue, dann %************************************************************************
% Gemeinsame Preamble fuer Skript (docu.tex), Folien (slides.tex) und
% Handouts/Uebungsblaetter (handouts.tex).
% Jede Root-Datei laedt ihre Dokumentklasse, setzt ggf. \slidestrue bzw.
% \handoutstrue, dann %************************************************************************
% Gemeinsame Preamble fuer Skript (docu.tex), Folien (slides.tex) und
% Handouts/Uebungsblaetter (handouts.tex).
% Jede Root-Datei laedt ihre Dokumentklasse, setzt ggf. \slidestrue bzw.
% \handoutstrue, dann \input{preamble}, dann hyperref, dann die Geometrie.
% Aufbau und RWU-Design wie im Projekt Embedded Computing.
%************************************************************************
\usepackage{iftex}
\ifPDFTeX
  % Rueckfallebene ohne RWU-Schrift: make PDFTEX=1 (oder pdflatex)
  \usepackage[utf8]{inputenc}
  \usepackage[T1]{fontenc}
  \usepackage{lmodern}
\else
  % RWU-Hausschrift Barlow Semi Condensed, liegt in fonts/ (SIL OFL)
  \usepackage{fontspec}
  \defaultfontfeatures{Ligatures=TeX}
  \newcommand{\rwufontsetup}[1]{#1{BarlowSemiCondensed}[
      Path=fonts/, Extension=.otf,
      UprightFont=*-Regular, ItalicFont=*-Italic,
      BoldFont=*-SemiBold, BoldItalicFont=*-SemiBoldItalic]}
  \rwufontsetup\setmainfont
  \rwufontsetup\setsansfont
  \setmonofont{Latin Modern Mono}
\fi
\usepackage[ngerman]{babel}
\usepackage[autostyle=true,german=quotes]{csquotes}
\usepackage[table]{xcolor}
\usepackage{graphicx,amsmath,amssymb,latexsym,longtable,url,version,ifthen}
\usepackage{array,tabularx,booktabs,multicol}

% RWU-Hausfarben (wie rwudefs.sty)
\definecolor{rwuviolet}{RGB}{102,56,182}
\definecolor{rwuvioletlight}{RGB}{155,125,212}
\definecolor{rwucyan}{RGB}{0,169,206}
\definecolor{rwucyanlight}{RGB}{5,195,222}
\definecolor{rwucyan40}{RGB}{153,221,253}
\colorlet{rwucyandark}{rwucyan!75!black}   % Cyan fuer Text auf Weiss
\colorlet{rwugray}{black!55}
% zieht den ersten Spiegelstrich unter ein Themenlabel:
% \item[{\rot\bf Thema}]~\\[-\lblskip]
\newcommand{\lblskip}{1.1em}
\newcommand{\rwulogo}[1][height=1.3cm]{\includegraphics[#1]{rwu_logo_hor_lila-cyan_cmyk}}
\usepackage{chngcntr}
\usepackage{verbatim}
\usepackage{tikz}
\usetikzlibrary{arrows.meta,positioning,calc,shapes.geometric,decorations.pathreplacing}
\usepackage[most]{tcolorbox}

% Literatur: biblatex/biber, Stil wie im Buch (authoryear)
\usepackage[backend=biber,style=authoryear,maxcitenames=2,uniquename=false,uniquelist=false,backref=false]{biblatex}
\addbibresource{refs.bib}
\setcounter{biburlnumpenalty}{9000}\setcounter{biburlucpenalty}{9000}\setcounter{biburllcpenalty}{9000}
, dann hyperref, dann die Geometrie.
% Aufbau und RWU-Design wie im Projekt Embedded Computing.
%************************************************************************
\usepackage{iftex}
\ifPDFTeX
  % Rueckfallebene ohne RWU-Schrift: make PDFTEX=1 (oder pdflatex)
  \usepackage[utf8]{inputenc}
  \usepackage[T1]{fontenc}
  \usepackage{lmodern}
\else
  % RWU-Hausschrift Barlow Semi Condensed, liegt in fonts/ (SIL OFL)
  \usepackage{fontspec}
  \defaultfontfeatures{Ligatures=TeX}
  \newcommand{\rwufontsetup}[1]{#1{BarlowSemiCondensed}[
      Path=fonts/, Extension=.otf,
      UprightFont=*-Regular, ItalicFont=*-Italic,
      BoldFont=*-SemiBold, BoldItalicFont=*-SemiBoldItalic]}
  \rwufontsetup\setmainfont
  \rwufontsetup\setsansfont
  \setmonofont{Latin Modern Mono}
\fi
\usepackage[ngerman]{babel}
\usepackage[autostyle=true,german=quotes]{csquotes}
\usepackage[table]{xcolor}
\usepackage{graphicx,amsmath,amssymb,latexsym,longtable,url,version,ifthen}
\usepackage{array,tabularx,booktabs,multicol}

% RWU-Hausfarben (wie rwudefs.sty)
\definecolor{rwuviolet}{RGB}{102,56,182}
\definecolor{rwuvioletlight}{RGB}{155,125,212}
\definecolor{rwucyan}{RGB}{0,169,206}
\definecolor{rwucyanlight}{RGB}{5,195,222}
\definecolor{rwucyan40}{RGB}{153,221,253}
\colorlet{rwucyandark}{rwucyan!75!black}   % Cyan fuer Text auf Weiss
\colorlet{rwugray}{black!55}
% zieht den ersten Spiegelstrich unter ein Themenlabel:
% \item[{\rot\bf Thema}]~\\[-\lblskip]
\newcommand{\lblskip}{1.1em}
\newcommand{\rwulogo}[1][height=1.3cm]{\includegraphics[#1]{rwu_logo_hor_lila-cyan_cmyk}}
\usepackage{chngcntr}
\usepackage{verbatim}
\usepackage{tikz}
\usetikzlibrary{arrows.meta,positioning,calc,shapes.geometric,decorations.pathreplacing}
\usepackage[most]{tcolorbox}

% Literatur: biblatex/biber, Stil wie im Buch (authoryear)
\usepackage[backend=biber,style=authoryear,maxcitenames=2,uniquename=false,uniquelist=false,backref=false]{biblatex}
\addbibresource{refs.bib}
\setcounter{biburlnumpenalty}{9000}\setcounter{biburlucpenalty}{9000}\setcounter{biburllcpenalty}{9000}
, dann hyperref, dann die Geometrie.
% Aufbau und RWU-Design wie im Projekt Embedded Computing.
%************************************************************************
\usepackage{iftex}
\ifPDFTeX
  % Rueckfallebene ohne RWU-Schrift: make PDFTEX=1 (oder pdflatex)
  \usepackage[utf8]{inputenc}
  \usepackage[T1]{fontenc}
  \usepackage{lmodern}
\else
  % RWU-Hausschrift Barlow Semi Condensed, liegt in fonts/ (SIL OFL)
  \usepackage{fontspec}
  \defaultfontfeatures{Ligatures=TeX}
  \newcommand{\rwufontsetup}[1]{#1{BarlowSemiCondensed}[
      Path=fonts/, Extension=.otf,
      UprightFont=*-Regular, ItalicFont=*-Italic,
      BoldFont=*-SemiBold, BoldItalicFont=*-SemiBoldItalic]}
  \rwufontsetup\setmainfont
  \rwufontsetup\setsansfont
  \setmonofont{Latin Modern Mono}
\fi
\usepackage[ngerman]{babel}
\usepackage[autostyle=true,german=quotes]{csquotes}
\usepackage[table]{xcolor}
\usepackage{graphicx,amsmath,amssymb,latexsym,longtable,url,version,ifthen}
\usepackage{array,tabularx,booktabs,multicol}

% RWU-Hausfarben (wie rwudefs.sty)
\definecolor{rwuviolet}{RGB}{102,56,182}
\definecolor{rwuvioletlight}{RGB}{155,125,212}
\definecolor{rwucyan}{RGB}{0,169,206}
\definecolor{rwucyanlight}{RGB}{5,195,222}
\definecolor{rwucyan40}{RGB}{153,221,253}
\colorlet{rwucyandark}{rwucyan!75!black}   % Cyan fuer Text auf Weiss
\colorlet{rwugray}{black!55}
% zieht den ersten Spiegelstrich unter ein Themenlabel:
% \item[{\rot\bf Thema}]~\\[-\lblskip]
\newcommand{\lblskip}{1.1em}
\newcommand{\rwulogo}[1][height=1.3cm]{\includegraphics[#1]{rwu_logo_hor_lila-cyan_cmyk}}
\usepackage{chngcntr}
\usepackage{verbatim}
\usepackage{tikz}
\usetikzlibrary{arrows.meta,positioning,calc,shapes.geometric,decorations.pathreplacing}
\usepackage[most]{tcolorbox}

% Literatur: biblatex/biber, Stil wie im Buch (authoryear)
\usepackage[backend=biber,style=authoryear,maxcitenames=2,uniquename=false,uniquelist=false,backref=false]{biblatex}
\addbibresource{refs.bib}
\setcounter{biburlnumpenalty}{9000}\setcounter{biburlucpenalty}{9000}\setcounter{biburllcpenalty}{9000}
, dann hyperref, dann die Geometrie.
% Aufbau und RWU-Design wie im Projekt Embedded Computing.
%************************************************************************
\usepackage{iftex}
\ifPDFTeX
  % Rueckfallebene ohne RWU-Schrift: make PDFTEX=1 (oder pdflatex)
  \usepackage[utf8]{inputenc}
  \usepackage[T1]{fontenc}
  \usepackage{lmodern}
\else
  % RWU-Hausschrift Barlow Semi Condensed, liegt in fonts/ (SIL OFL)
  \usepackage{fontspec}
  \defaultfontfeatures{Ligatures=TeX}
  \newcommand{\rwufontsetup}[1]{#1{BarlowSemiCondensed}[
      Path=fonts/, Extension=.otf,
      UprightFont=*-Regular, ItalicFont=*-Italic,
      BoldFont=*-SemiBold, BoldItalicFont=*-SemiBoldItalic]}
  \rwufontsetup\setmainfont
  \rwufontsetup\setsansfont
  \setmonofont{Latin Modern Mono}
\fi
\usepackage[ngerman]{babel}
\usepackage[autostyle=true,german=quotes]{csquotes}
\usepackage[table]{xcolor}
\usepackage{graphicx,amsmath,amssymb,latexsym,longtable,url,version,ifthen}
\usepackage{array,tabularx,booktabs,multicol}

% RWU-Hausfarben (wie rwudefs.sty)
\definecolor{rwuviolet}{RGB}{102,56,182}
\definecolor{rwuvioletlight}{RGB}{155,125,212}
\definecolor{rwucyan}{RGB}{0,169,206}
\definecolor{rwucyanlight}{RGB}{5,195,222}
\definecolor{rwucyan40}{RGB}{153,221,253}
\colorlet{rwucyandark}{rwucyan!75!black}   % Cyan fuer Text auf Weiss
\colorlet{rwugray}{black!55}
% zieht den ersten Spiegelstrich unter ein Themenlabel:
% \item[{\rot\bf Thema}]~\\[-\lblskip]
\newcommand{\lblskip}{1.1em}
\newcommand{\rwulogo}[1][height=1.3cm]{\includegraphics[#1]{rwu_logo_hor_lila-cyan_cmyk}}
\usepackage{chngcntr}
\usepackage{verbatim}
\usepackage{tikz}
\usetikzlibrary{arrows.meta,positioning,calc,shapes.geometric,decorations.pathreplacing}
\usepackage[most]{tcolorbox}

% Literatur: biblatex/biber, Stil wie im Buch (authoryear)
\usepackage[backend=biber,style=authoryear,maxcitenames=2,uniquename=false,uniquelist=false,backref=false]{biblatex}
\addbibresource{refs.bib}
\setcounter{biburlnumpenalty}{9000}\setcounter{biburlucpenalty}{9000}\setcounter{biburllcpenalty}{9000}
